\begin{name}
	{\tenchude}
	{TOÁN 12 - ÔN TẬP CHƯƠNG V}
	{LỚP TOÁN THẦY PHÁT}
	{Thời gian: 90 phút - Không kể thời gian phát đề}
\end{name}
\caulc
\Opensolutionfile{ans}[ans/ans-2-C5-OTC-De2-LC]
\begin{ex}%[2H5N3-3]
	Trong hệ trục tọa độ $Oxyz$, phương trình mặt cầu tâm $I(2;1;-2)$ bán kính $R=2$ là
	\choice
	{$(x-2)^2+(y-1)^2+(z-2)^2=2^2$}
	{\True $(x-2)^2+(y-1)^2+(z+2)^2=2^2$}
	{$(x+2)^2+(y+1)^2+(z-2)^2=4$}
	{$(x-2)^2+(y-1)^2+(z+2)^2=2$}
	\loigiai{
		Phương trình mặt cầu tâm $I(2;1;-2)$ bán kính $R=2$ có dạng $$(x-2)^2+(y-1)^2+(z+2)^2=2^2.$$}
\end{ex}

\begin{ex}%[2H5H3-3]
	Trong không gian với hệ tọa độ $Oxyz$, viết phương trình mặt cầu có tâm $I(1;-4;3)$ và đi qua điểm $A(5;-3;2)$.
	\choice
	{$(x-1)^2+(y-4)^2+(z-3)^2=18$}
	{$(x-1)^2+(y-4)^2+(z-3)^2=16$}
	{$(x-1)^2+(y+4)^2+(z-3)^2=16$}
	{\True $(x-1)^2+(y+4)^2+(z-3)^2=18$}
	\loigiai{
		Mặt cầu có tâm $I(1;-4;3)$ và đi qua điểm $A(5;-3;2)$ nên có bán kính $$R=IA=\sqrt{(5-1)^2+(-3+4)^2+(2-3)^2}=3\sqrt{2}.$$
		Vậy phương trình mặt cầu cần tìm là $(x-1)^2+(y+4)^2+(z-3)^2=18$.}
\end{ex}
\begin{ex}%[2H5N2-2]
	Trong không gian với hệ tọa độ $Oxyz$, vectơ nào dưới đây là một vectơ chỉ phương của đường thẳng $d\colon \heva{& x=1+t\\& y=4\\& z=3-2t.}$ 
	\choice
	{$\overrightarrow{u}=(1;4;3)$}
	{$\overrightarrow{u}=(1;4;-2)$}
	{\True $\overrightarrow{u}=(1;0;-2)$}
	{$\overrightarrow{u}=(1;0;2)$}
	\loigiai{
		Từ phương trình tham số của đường thẳng $d$, ta suy ra một vectơ chỉ phương của đường thẳng $d$ là $\overrightarrow{u}=(1;0;-2)$.}
\end{ex}

\begin{ex}%[2H5N2-3]
	Trong không gian với hệ tọa độ $Oxyz$, cho đường thẳng $\Delta$ đi qua điểm $M(2;0;-1)$ và có một vectơ chỉ phương $\overrightarrow{a}=(4;-6;2)$. Phương trình tham số của $\Delta$ là
	\choice
	{$\heva{& x=-2+4t\\& y=6t\\& z=1+2t}$}
	{\True $\heva{& x=2+2t\\& y=-3t\\& z=-1+t}$}
	{$\heva{& x=4+2t\\& y=-6\\& z=2+t}$}
	{$\heva{& x=-2+2t\\& y=3t\\& z=1+t}$}
	\loigiai{
		Ta có $\overrightarrow{a}=(4;-6;2)=2(2;-3;1)$.\\
		Do đó đường thẳng $\Delta$ có một vectơ chỉ phương là $\overrightarrow{u}=(2;-3;1)$.\\
		Vậy phương trình tham số của $\Delta$ đi qua $M(2;0;-1)$ và có một vectơ chỉ phương là $\overrightarrow{u}=(2;-3;1)$ là $\heva{& x=2+2t\\& y=-3t\\& z=-1+t.}$}
\end{ex}
\begin{ex}%[2H5H2-3]
	Trong không gian với hệ trục $Oxyz$, cho điểm $M(1;-3;2)$ và mặt phẳng   $(P)\colon x-3y+2z-1=0$. Tìm phương trình đường thẳng $d$ qua $M$ và vuông góc với $(P)$.
	\choice
	{$\dfrac{x+1}{1}=\dfrac{y-3}{-3}=\dfrac{z+2}{2}$}
	{\True $\dfrac{x-1}{1}=\dfrac{y+3}{-3}=\dfrac{z-2}{2}$ }
	{$\dfrac{x}{1}=\dfrac{y}{-3}=\dfrac{z}{2}$} {$\dfrac{x+1}{1}=\dfrac{y+3}{-3}=\dfrac{z-2}{2}$}
	\loigiai{
		Mặt phẳng $(P)$ có vectơ pháp tuyến là $\overrightarrow{n}=(1;-3;2)$.\\
		Vì $d$ vuông góc với $(P)$ nên $d$ nhận $\overrightarrow{n}=(1;-3;2)$ là vectơ chỉ phương.\\
		Đường thẳng $d$ qua $M$ và nhận $\overrightarrow{n}=(1;-3;2)$ là vectơ chỉ phương có phương trình $\dfrac{x-1}{1}=\dfrac{y+3}{-3}=\dfrac{z-2}{2}$.}
\end{ex}
		
\begin{ex}%[2H5H2-5] 
	Trong không gian với hệ tọa độ $Oxyz$, khoảng cách giữa đường thẳng   $d\colon \dfrac{x-1}{2}=\dfrac{y-3}{2}=\dfrac{z-2}{1}$ và mặt phẳng $(P)\colon x-2y+2z+4=0$ bằng bao nhiêu?
	\choice
	{\True $1$}
	{$0$}
	{$3$}
	{$2$}
	\loigiai{
	Mặt phẳng $(P)$ có vectơ pháp tuyến là $\overrightarrow{n}=(1;-2;2)$.\\
	Đường thẳng $d$ có vectơ chỉ phương là $\overrightarrow{u}=(2;2;1)$.\\
	Ta có $\overrightarrow{n}\cdot\overrightarrow{u}=1\cdot 2+(-2)\cdot 2+2\cdot 1=0$ nên $\overrightarrow{n}\perp\overrightarrow{u}$.\\
	Suy ra đường thẳng $d$ song song với mặt phẳng $(P)$.\\
	Chọn $M(1;3;2)\in d$.\\
	$\mathrm{d}(d;(P))=\mathrm{d}(M;(P))=\dfrac{\left| 1-6+4+4\right|}{\sqrt{1^2+(-2)^2+2^2}}=1$.}
\end{ex}
	
\begin{ex}%[2H5N1-2]
	Trong không gian với hệ tọa độ $Oxyz$, cho phương trình tổng quát của mặt phẳng $(P)\colon 2x-6y-8z+1=0$. Một vectơ pháp tuyến của mặt phẳng $(P)$ có tọa độ là
	\choice
	{$(-1;-3;4)$}
	{$(1;3;4)$}
	{\True $(1;-3;-4)$}
	{$(1;-3;4)$}
	\loigiai{
	Phương trình tổng quát của mặt phẳng $(P)\colon 2x-6y-8z+1=0$ nên một vectơ pháp tuyến của mặt phẳng $(P)$ có tọa độ là $(2;-6;-8)$ hay $(1;-3;4)$.}
\end{ex}
\begin{ex}%[2H5H1-3]
	Trong không gian với hệ trục tọa độ $Oxyz$, phương trình nào dưới đây là phương trình mặt phẳng đi qua $M(1;-1;2)$ và vuông góc với đường thẳng $\Delta\colon\dfrac{x+1}{2}-\dfrac{y-2}{-1}=\dfrac{z}{3}$?
	\choice
	{$2x-y+3z+9=0$}
	{$2x+y+3z-9=0$}
	{\True $2x-y+3z-9=0$}
	{$2x-y+3z-6$}
	\loigiai{
	Mặt phẳng $(P)$ vuông góc với $\Delta$ nên $(P)$ nhận vectơ chỉ phương của $\Delta$ là $\overrightarrow{u}_{\Delta}=(2;-1;3)$ làm vectơ pháp tuyến\\
		$\Rightarrow $Phương trình mặt phẳng $(P)$ là $2(x-1)-1(y+1)+3(z-2)=0$ hay $2x-y+3z-9=0$.}
\end{ex}

\begin{ex}%[2H5H1-3]
	Trong không gian với hệ trục tọa độ $Oxyz$, mặt phẳng đi qua điểm $A(0;1;0)$ và chứa đường thẳng $(\Delta)\colon\dfrac{x-2}{1}=\dfrac{y-1}{-1}=\dfrac{z-3}{1}$ có phương trình là
	\choice
	{$x-y+z+1=0$}
	{$3x-y+2z+1=0$}
	{$x+y+z-1=0$}
	{\True $3x+y-2z-1=0$}
	\loigiai{
	Ta lấy điểm $M(2;1;3)\in (\Delta)\Rightarrow\heva{&\overrightarrow{AM}=(2;0;3)\\&
	\text{vectơ chỉ phương} \overrightarrow{u}_{(\Delta)}=(1;-1;1)}$  $\Rightarrow\overrightarrow{n}=\left[\overrightarrow{AM},\overrightarrow{u}_{\left(\Delta\right)}\right]=(3;1;-2)$.\\
	Mặt phẳng cần tìm qua $A(0;1;0)$ và nhận $\overrightarrow{n}=(3;1;-2)$ làm vectơ pháp tuyến có phương trình là $3(x-0)+1(y-1)-2(z-0)=0\Leftrightarrow 3x+y-2z-1=0$.}
\end{ex}

\begin{ex}%[2H5H1-3]
	Trong không gian với hệ trục tọa độ $Oxyz$, phương trình của mặt phẳng $(P)$ đi qua điểm $B(2;1;-3)$, đồng thời vuông góc với hai mặt phẳng $(Q)\colon x+y+3z=0$, $(R)\colon 2x-y+z=0$ là
	\choice
	{$4x+5y-3z+22=0$}
	{$4x-5y-3z-12=0$}
	{$2x+y-3z-14=0$}
	{\True $4x+5y-3z-22=0$}
	\loigiai{
	Mặt phẳng $(Q)\colon x+y+3z=0$, $(R)\colon 2x-y+z=0$ có các vectơ pháp tuyến lần lượt là $\overrightarrow{n}_1=(1;1;3)$ và $\overrightarrow{n}_2=(2;-1;1)$.\\
	Vì $(P)$ vuông góc với hai mặt phẳng $(Q)$, $(R)$ nên $(P)$ có vectơ pháp tuyến là  $\overrightarrow{n}=\left[\overrightarrow{n}_1,\overrightarrow{n}_2\right]=(4;5;-3)$.\\
	Ta lại có $(P)$ đi qua điểm $B(2;1;-3)$ nên $$(P)\colon 4(x-2)+5(y-1)-3(z+3)=0\Leftrightarrow 4x+5y-3z-22=0.$$}
\end{ex}

\begin{ex}%[2H5H1-3]
	Trong không gian với hệ trục tọa độ $Oxyz$, gọi ba điểm $A$, $B$, $C$ lần lượt là hình chiếu vuông góc của điểm $M(1;2;3)$ lên các trục tọa độ. Viết phương trình mặt phẳng $(ABC)$. 
	\choice
	{$\dfrac{1}{x}+\dfrac{2}{y}+\dfrac{3}{z}=1$}
	{\True $\dfrac{x}{1}+\dfrac{y}{2}+\dfrac{z}{3}=1$}
	{$\dfrac{1}{x}+\dfrac{2}{y}+\dfrac{3}{z}=0$}
	{$\dfrac{x}{1}+\dfrac{y}{2}+\dfrac{z}{3}=0$}
	\loigiai{
		Gọi $A$, $B$, $C$ lần lượt là hình chiếu vuông góc của điểm $M(1;2;3)$ lên $Ox$, $Oy$, $Oz$.\\
		Suy ra $A(1;0;0)$, $B(0;2;0)$, $C(0;0;3)$.\\
		Vậy phương trình mặt phẳng $(ABC)$ theo đoạn chắn là $\dfrac{x}{1}+\dfrac{y}{2}+\dfrac{z}{3}=1$.}
\end{ex}

\begin{ex}%[2H5V2-4]
	Trong không gian với hệ tọa độ $Oxyz$, cho hai đường thẳng $d_1\colon\dfrac{x-1}{2}=\dfrac{y}{1}=\dfrac{z}{3}$, $d_2\colon\heva{& x=1+t\\& y=2+t\\& z=m}$. Gọi $S$ là tập tất cả các số $m$ sao cho $d_1$ và $d_2$ chéo nhau và khoảng cách giữa chúng bằng $\dfrac{5}{\sqrt{19}}$. Tính tổng các phần tử của $S$.
	\choice
	{$-11$}
	{$12$}
	{\True $-12$}
	{$11$}
	\loigiai{
	$d_1$ đi qua điểm $M(1;0;0)$, có vectơ chỉ phương $\vec{u}_1=(2;1;3)$.\\
	$d_2$ đi qua điểm $N(1;2;m)$, có vectơ chỉ phương $\vec{u}_2=(1;1;0)$.\\
	Ta có $\left[\vec{u}_1,\vec{u}_2\right]=(-3;3;1)$, $\overrightarrow{MN}=(0;2;m)$.\\
	$d_1$ và $d_2$ chéo nhau khi và chỉ khi $\left[\vec{u}_1,\vec{u}_2\right]\cdot\overrightarrow{MN}\ne 0\Leftrightarrow m\ne-6$.\\
	Mặt khác $\mathrm{d}\left(d_1,d_2\right)=\dfrac{5}{\sqrt{19}}$ $\Leftrightarrow\dfrac{\left|\left[\vec{u}_1,\vec{u}_2\right]\cdot\overrightarrow{MN}\right|}{\left|\left[\vec{u}_1,\vec{u}_2\right]\right|}=\dfrac{5}{\sqrt{19}}$ $\Leftrightarrow\dfrac{\left| m+6\right|}{\sqrt{19}}=\dfrac{5}{\sqrt{19}}$ $\Leftrightarrow\heva{&m=-1\\&m=-11.}$\\
	Khi đó tổng các phần tử của $m$ là $-12$.}
\end{ex}
\Closesolutionfile{ans}
\indapan{6}{ans/ans-2-C5-OTC-De2-LC}

\cauds
\Opensolutionfile{ans}[ans/ans-2-C5-OTC-De2-DS]
% \setcounter{ex}{0}
\begin{ex}%[2H5V3-3]
	Trong không gian $Oxyz$, cho mặt cầu $(S)\colon x^2+(y-1)^2+(z+2)^2=9$. 
	\choiceTF
	{\True Tâm của mặt cầu $(S)$ là điểm $I(0;1;-2)$}
	{Điểm $A(3;1;0)$ nằm trên mặt cầu $(S)$}
	{Tâm mặt cầu $(S)$ cách mặt phẳng $(Oyz)$ một khoảng bằng $2$}
	{Khi mặt phảng $(P)\colon x+my-z-2=0$ cắt mặt cầu $(S)$ theo đường tròn có bán kính bằng $3$ thì giá trị $m=2$}
	\loigiai{
		\begin{itemchoice}
			\itemch \textbf{Đúng}. Vì $(S)\colon x^2+(y-1)^2+(z+2)^2=9\Rightarrow I(0; 1;-2)$ là tâm của $(S)$.
			\itemch \textbf{Sai}. Vì  $3^2+(1-1)^2+(0+2)^2\ne 9\Rightarrow A(3;1; 0)\notin(S)$.
			\itemch \textbf{Sai}. Vì $I(0;1;-2)\Rightarrow I\in(Oyz)\Rightarrow \mathrm{d}\left(I;(Oyz)\right)=0$.
			\itemch \textbf{Sai}. Vì  $(S)$ có bán kính $R=3$ nên để mặt phẳng $(P)\colon x+my-z-2=0$ cắt mặt cầu $(S)$ theo đường tròn có bán kính bằng $3$ thì tâm $I(0;1;-2)\in (P)$\\ $\Rightarrow 0+ m\cdot 1-(-2)-2=0\Leftrightarrow m=0$.
	\end{itemchoice}}
\end{ex}
\begin{ex}%[2H5V3-3]
	Trong không gian $Oxyz$, cho đường thẳng $d$ đi qua gốc tọa độ $O$ có $\vec{v}=(1; 2;3)$ là vectơ chỉ phương và mặt cầu $(S)\colon x^2+y^2+(z-1)^2=9$.
	\choiceTF
	{\True Đường thẳng $d$ có phương trình chính tắc là $\dfrac{x}{1}=\dfrac{y}{2}=\dfrac{z}{3}$}
	{Đường thẳng $d$ đi qua $A(-1;2;3)$}
	{Đường thẳng $d$ và mặt cầu $(S)$ có một điểm chung duy nhất là $B(1;2;3)$}
	{\True  Khi mặt phẳng $(P)$ song song $d$ và trục $Ox$ thì $(P)$ nhận $\vec{a}=(0 ;-3;2)$ làm vectơ pháp tuyến}
	\loigiai{
		\begin{itemchoice}
			\itemch \textbf{Đúng}. Vì $O(0;0;0) \Rightarrow d$ có phương trình chính tắc là $\dfrac{x}{1}=\dfrac{y}{2}=\dfrac{z}{3}$.
			\itemch \textbf{Sai}. Vì $\dfrac{-1}{1}\ne \dfrac{2}{2}=\dfrac{3}{3} \Rightarrow A(-1;2;3) \notin d$.
			\itemch \textbf{Sai}. Vì  $d$ có phương trình tham số là $\heva{&x=t\\&y=2 t\\&z=3 t.}$\\
			Xét hệ phương trình
			\allowdisplaybreaks
			\begin{eqnarray*}
				\heva{&x=t\\&y=2t\\&z=3t\\&	x^2+y^2+(z-1)^2=9} 
					&\Leftrightarrow& \heva{&x=t\\&y=2t\\&z=3t\\& t^2+(2t)^2+(3t-1)^2=9}\\
					&\Leftrightarrow& \heva{&t=1\\& x=1\\&y=2\\&z=3} \text{ hoặc } \heva{&t=-\dfrac{4}{7} \\& x=-\dfrac{4}{7}\\&y=-\dfrac{8}{7}\\&z=-\dfrac{12}{7}.}
			\end{eqnarray*}
			Suy ra $ d$ và $(S)$ có $2$ điểm chung.
			\itemch \textbf{Đúng}. Vì trục $Ox$ có vectơ chỉ phương là $\vec{i}=(1;0;0)$ nên để mặt phẳng $(P)$ song song $d$ và trục $Ox$ thì $(P)$ nhận $\vec{n}=\left[\vec{v};\vec{i}\right]=(0;3;-2)$ làm vectơ pháp tuyến.\\
			Mà $\vec{a}=-\vec{n} \Rightarrow(P)$ cũng nhận $\vec{a}=(0;-3;2)$ làm vectơ pháp tuyến.
		\end{itemchoice}
}\end{ex}

\begin{ex}%[2H5V1-6]
	Trong không gian $Oxyz$, cho mặt phẳng $(P)$ đi qua $3$ điểm $A(1;0;0)$, $B(0;-2;0)$, $C(0;0;1)$. 
	\choiceTF
	{Mặt phẳng $(P)$ có phương trình là $2x-y+2z=1$}
	{\True Mặt phẳng $(P)$ đi qua $D(-1;2;3)$}
	{$\cos \left((P);(Oxy)\right)=\dfrac{\sqrt{2}}{2}$}
	{\True Hình chóp $O.ABC$ có chiều cao $h=\dfrac{2}{3}$}
	\loigiai{
		\begin{itemchoice}
			\itemch \textbf{Sai}. Vì $A(1;0;0)\in Ox$, $B(0;-2;0)\in Oy$, $C(0;0;1)\in Oz$ nên phương trình mp$(P)$ là $$\dfrac{x}{1}-\dfrac{y}{2}+\dfrac{z}{1}=1 \Leftrightarrow 2x-y+2z=2$$
			\itemch \textbf{Đúng}. Vì $2\cdot (-1)-2+2\cdot 3=2\Rightarrow D(-1;2;3) \in(P)$.
			\itemch \textbf{Sai}. Vì $(P)$ có vectơ pháp tuyến là $\vec{n}=(2;-1;2)$.
			$(Oxy)$ có vectơ pháp tuyến là $\vec{j}=(0;0;1)$ nên
			$$\cos \left((P);(Oxy)\right)=\dfrac{\left|2\cdot 0+(-1)\cdot 0+2 \cdot 1\right|}{\sqrt{2^{2}+(-1)^{2}+2^{2}} \sqrt{0^{2}+0^{2}+1^{2}}}=\dfrac{2}{3}.$$
			\itemch \textbf{Đúng}. Vì ta có  $h=\mathrm{d}(O;(P))=\dfrac{\left|2\cdot 0-0+2\cdot 0-2\right|}{\sqrt{2^2+(-1)^2+2^{2}}}=\dfrac{2}{3}$.
		\end{itemchoice}
}\end{ex}

\begin{ex}%[2H5V3-3] 
	Trong không gian $Oxyz$, cho mặt phẳng $(P)\colon x+3y-z-2=0$ và mặt cầu
	$(S)\colon x^2+y^2+z^2-2x+4y-11=0$. 
	\choiceTF
	{\True Mặt phẳng $(P)$ nhận $\vec{n}=(1;3;-1)$ làm vectơ pháp tuyến}
	{\True  Mặt cầu $(S)$ có bán kính $R=4$}
	{\True Đường thẳng $d$ đi qua tâm $I$ của mặt cầu $(S)$ và vuông góc với mặt phẳng $(P)\colon \dfrac{x-2}{1}=\dfrac{y-1}{3}=\dfrac{z+1}{-1}$}
	{\True Mặt phẳng $(P)$ và mặt cầu $(S)$ có vô số điểm chung}
	\loigiai{
		\begin{itemchoice}
			\itemch \textbf{Đúng}.  Vì mặt phẳng $(P)\colon  x+3 y-z-2=0 \Rightarrow$ mặt phẳng $(P)$ nhận $\vec{n}=(1;3;-1)$ làm vectơ pháp tuyến.
			\itemch \textbf{Đúng}.  Vì 
			\allowdisplaybreaks
			\begin{eqnarray*}
				&&(S)\colon  x^2+y^2+z^2-2 x+4 y-11=0\\&\Rightarrow& a=\dfrac{-2}{-2}=1,\, b=\dfrac{4}{-2}=-2,\, c=\dfrac{0}{-2}=0,\, d=-11 \\&\Rightarrow& R=\sqrt{1^{2}+(-2)^{2}+0^{2}-(-11)}=4.
			\end{eqnarray*}
			\itemch \textbf{Đúng}.  Vì  $(S)$ có tâm $I(1;-2;0)$, mặt phẳng $(P)$ có vectơ pháp tuyến $\vec{n}=(1;3;-1)$.\\
			Đường thẳng $d$ đi qua tâm $I$ của mặt cầu $(S)$ và vuông góc với mặt phẳng $(P)$ nên $d$ nhận $\vec{n}=(1;3;-1)$ làm vectơ chỉ phương $\Rightarrow d\colon \dfrac{x-1}{1}=\dfrac{y+2}{3}=\dfrac{z}{-1}$.\\
			Ta thấy $E(2;1;-1)\in d$, nên phương trình $d$ có thể viết lại là $\dfrac{x-2}{1}=\dfrac{y-1}{3}=\dfrac{z+1}{-1}$.
			\itemch \textbf{Đúng}.  Vì  $\mathrm{d}(I;(P))=\dfrac{\left|1\cdot 1+3 \cdot(-2)-1 \cdot 0-2\right| }{\sqrt{1^{2}+3^{2}+(-1)^{2}}}=\dfrac{7}{\sqrt{11}}<R=4$ nên mặt phẳng $(P)$ cắt mặt cầu $(S)$ theo thiết diện là một đường tròn $\Rightarrow$ mặt phẳng $(P)$ và mặt cầu $(S)$ có vô số điểm chung.
		\end{itemchoice}
}
\end{ex}
\Closesolutionfile{ans}
\indapan{2}{ans/ans-2-C5-OTC-De2-DS}

\caukq
\Opensolutionfile{ans}[ans/ans-2-C5-OTC-De2-KQ]
\begin{ex}%[2H5H1-5]  
	Trong không gian với hệ trục tọa độ $Oxyz$, cho mặt phẳng  $(\alpha)\colon 6x-3y-2z-6=0$. Khoảng cách từ $O$ đến mặt phẳng $(\alpha)$ bằng bao nhiêu? (Kết quả làm tròn đến hàng phần trăm).
	
	\shortans{0,86}
	\loigiai{
		Ta có $\mathrm{d}(O,(\alpha))=\dfrac{6}{\sqrt{36+9+4}}=\dfrac{6}{7}\approx 0{,}86$.
}
\end{ex}
\begin{ex}%[2H5C3-4]
Trong không gian $Oxyz$, tỉ lệ dài trên các trục là $10$ km tính cho một đơn vị tỉ lệ trên mỗi trục, cho bốn vệ tinh có toạ độ lần lượt là $A(-1;2;-1)$, $B(1;4;0)$, $C(3;0;9)$, $D(7;10;-1)$. Đang tiến hành theo dõi vật thể $ M $ (coi là một chất điểm). Ở một thời điểm cả 4 vệ tinh bắn tín hiệu về $M$ và đo được $MA = 6$, $MB = 3$, $MC = 10$, $MD = 6$. Ngay sau đó $10$ giây cả $4$ vệ tinh lại bắn tín hiệu về vật $ M $ và đo được $ MA = 3 $, $ MB = 4 $, $ MC = \sqrt{85} $, $ MD = \sqrt{137} $. Nếu coi như vật $ M $ chuyển động thẳng đều thì tốc độ của vật $ M $ bằng bao nhiêu km/s? (Kết quả làm tròn đến hai chữ số thập phân).
\shortans[oly]{$0{,}64$}
\loigiai
{
Đặt $M_1(x_1;y_1;z_1)$ là tọa độ vật thể $M$ ở thời điểm ban đầu. \\
Ta có
$\heva{&M_1A=6\\&M_1B=3\\&M_1C=10\\&M_1D=6}\Leftrightarrow\heva{&M_1A^2=36\\&M_1B^2=9\\&M_1C^2=100\\&M_1D^2=36}\Leftrightarrow$
$\heva{&(x_1 + 1)^2 + (y_1 - 2)^2 + (z_1 + 1)^2 = 36 \\&(x_1 - 1)^2 + (y_1 - 4)^2 + z_1^2 = 9\\&(x_1 - 3)^2 + y_1^2 + (z_1 - 9)^2 = 100 \\&(x_1 - 7)^2 + (y_1 - 10)^2 + (z_1 + 1)^2 = 36 }$\\	$\Leftrightarrow$ $\heva{&x_1^2+y_1^2+z_1^2+2x_1-4y_1+2z_1=30\quad (1)\\&x_1^2+y_1^2+z_1^2-2x_1-8y_1=-8 \quad (2)\\&x_1^2+y_1^2+z_1^2-6x_1-18z_1=10\quad (3)\\&x_1^2+y_1^2+z_1^2-14x_1-20y_1+2z_1=-114\quad (4)}\Rightarrow$ $\heva{&4x_1 + 4y_1 + 2z_1 = 38  \quad \text{(lấy(1)-(2))} \\&8x_1-4y_1+20z_1=20\quad \text{(lấy(1)-(3))}\\&16x_1 + 16y_1 = 144\quad \text{(lấy(1)-(4))}}$\\
Giải hệ phương trình ta được  $\heva{&x_1=3\\&y_1=6\\&z_1=1}$. Vậy $M_1(3;6;1)$.
Tương tự đặt $M_2(x_2;y_2;z_2)$ là tọa độ vật thể $M$ ở thời điểm $t=10 $ giây. \\
Ta có
$\heva{&M_2A=3\\&M_2B=4\\&M_2C=\sqrt{85}\\&M_2D=\sqrt{137}}$ $\Leftrightarrow$  	$\heva{&M_2A^2=9\\&M_2B^2=16\\&M_2C^2=85\\&M_2D^2=137}\Leftrightarrow$
$\heva{&(x_2 + 1)^2 + (y_2 - 2)^2 + (z_2 + 1)^2 = 9 \\&(x_2 - 1)^2 + (y_2 - 4)^2 + z_2^2 = 16\\&(x_2 - 3)^2 + y_2^2 + (z_2 - 9)^2 = 85 \\&(x_2 - 7)^2 + (y_2 - 10)^2 + (z_2 + 1)^2 = 137 }$\\	$\Leftrightarrow$ $\heva{&x_2^2+y_2^2+z_2^2+2x_2-4y_2+2z_2=3\quad (5)\\&x_2^2+y_2^2+z_2^2-2x_2-8y_2=-1 \quad (6)\\&x_2^2+y_2^2+z_2^2-6x_2-18z_2=-5\quad (7)\\&x_2^2+y_2^2+z_2^2-14x_2-20y_2+2z_2=-13\quad (8)}\Rightarrow$ $\heva{&4x_2 + 4y_2 + 2z_2 = 4  \quad \text{(lấy(5)-(6))} \\&8x_2-4y_2+20z_2=8\quad \text{(lấy(5)-(7))}\\&16x_2 + 16y_2 = 16\quad \text{(lấy(5)-(8))}}$\\
Giải hệ phương trình ta được $\heva{&x_2=1\\&y_2=0\\&z_2=0}$. Vậy $M_2(1;0;0)$. \\
Quãng đường vật thể $M$ di chuyển là $M_1M_2=\sqrt{41}$.\\ Tốc độ của vật thể $M$ là $v=\dfrac{M_1M_2}{t}=\dfrac{\sqrt{41}}{10}\approx 0{,}64$ km/s.
}
\end{ex}

\begin{ex}%[2H5V3-2]
	Trong không gian $Oxyz$, mặt cầu $(S)$ có tâm thuộc đường thẳng $d\colon \dfrac{x-2}{1}=\dfrac{y-1}{3}=\dfrac{z+1}{-1}$ và đi qua $2$ điểm $A(2;1;-1)$, $B(0 ;1;1)$ có đường kính bằng bao nhiêu? (kết quả làm tròn đến hàng phần trăm).
	
	\shortans{6,63}
	\loigiai{
		$(S)$ có tâm $I\in d\colon \dfrac{x-2}{1}=\dfrac{y-1}{3}=\dfrac{z+1}{-1} \Rightarrow I(t+2;3t+1;-t-1)$.\\
		$(S)$ đi qua $2$ điểm $A(2;1;-1)$, $B(0;1;1)$ nên ta có
		\allowdisplaybreaks
		\begin{eqnarray*}
			IA=IB=R&\Rightarrow& t^2+(3t)^2+(-t)^2=(t+2)^2+(3t)^2+(-t-2)^2 \Leftrightarrow t=-1
			\\&\Rightarrow& I(1;-2;0) \Rightarrow d=2R=2IA=2\sqrt{11}\approx 6{,}63.
		\end{eqnarray*}
}
\end{ex}
\begin{ex}%[2H5H1-3] 
	Trong không gian $Oxyz$, cho mặt phẳng $(P)\colon x+3y-z-2=0$,$(Q)\colon 2x-4 y-z-3=0$ và $A(1;1;-1)$. Mặt phẳng $(R)$ đi qua $A$ vuông góc với cả hai mặt phẳng $(P)$ và $(Q)$ có phương trình $ax+by+cz+1=0$. Tính giá trị biểu thức $T=a-2b+3c$.
	
	\shortans{17,5}
	\loigiai{
		Do $(R)$ vuông góc với cả hai mặt phẳng $(P)$ và $(Q)$ nên $(R)$ có vectơ pháp tuyến  $\vec{n}=\left[\vec{n}_{Q};\vec{n}_{P}\right]=(7;1;10)$.\\
		$(R)$ đi qua điểm $A(1;1;-1) \Rightarrow(R)\colon 7x+y+10z+2=0$
		$$\Rightarrow a=3{,}5;\, b=0{,}5;\, c=5 \Rightarrow T=17{,}5.$$
}
\end{ex}

\begin{ex}%[2H2V2-2] 
	Trong không gian với hệ tọa độ $Oxyz$, cho hai điểm $A(1;2;-2)$ và $B\left(\dfrac{8}{3};\dfrac{4}{3};\dfrac{8}{3}\right)$. Biết $I(a;b;c)$ là tâm của đường tròn nội tiếp tam giác $OAB$. Tính giá trị $a-b+c$.
	
	\shortans{0}
	\loigiai{
	\begin{center}
	\begin{tikzpicture}[scale=.7, >=stealth]
\coordinate[label=above left:$O$] (O) at (-2,0);
\coordinate[label=below left:$A$] (A) at (-3,-4);
\coordinate[label=below right:$B$] (B) at (5,-4);
\pgfmathsetmacro\a{(sqrt(65))/(sqrt(17))} 
\coordinate [label=below:$D$](D) at ($(A)! 1/(1+\a) ! (B)$);
\pgfmathsetmacro\b{(sqrt(17))/8} 
\coordinate (E) at ($(B)! 1/(1+\b) ! (O)$);
\pgfmathsetmacro\c{(sqrt(41))/8} 
\coordinate (F) at ($(A)! 1/(1+\c) ! (O)$);
\coordinate(U) at ($(O)-(D)+(O)-(D)+(O)-(D)+(B)$);
\coordinate(V) at ($(A)! 2.3 ! (O)$);
%Tâm I
\draw[name path=done,color=white] (O)--(D);
\draw[name path=dtwo,color=white] (A)--(E);
\path [name intersections={of=done and dtwo,by=I}];
\draw (I) node[below left]{$I$};
\draw[->] (O)--(I);
\draw[->] (I)--(D);
\draw[->] (O)--(A);
\draw[->] (O)--(B);
\draw[->] (D)--(A);
\draw[->] (D)--(B);
\draw (A)--(I);
%Vẽ kí hiệu góc
\draw pic[draw, angle radius=5mm]{angle=A--O--D};	
\draw pic[draw, angle radius=6mm]{angle=D--O--B};	
\draw  (-0.9695760613404443,-2.4146928936958933) circle (1.585307106304106cm);
\foreach \diem in {A,B,O,D,I}	\fill (\diem)circle(1pt);		
\end{tikzpicture}
	\end{center}
		Ta có $\vec{OA}=(1;2;-2)$, $\vec{OB}=\left(\dfrac{8}{3};\dfrac{4}{3}; \dfrac{8}{3}\right)$, do đó $OA=3$, $OB=4$.\\
		Gọi $D$ là chân đường phân giác trong kẻ từ $O$.\\
		Ta có $\dfrac{AD}{DB}=\dfrac{OA}{OB}=\dfrac{3}{4}$, suy ra $\vec{DA}=-\dfrac{3}{4}\vec{DB} \Leftrightarrow 4\vec{DA}+3\vec{DB}=\vec{0} \Rightarrow \vec{OD}=\dfrac{4\cdot \vec{OA}+3 \cdot\vec{OB}}{7}$. Do đó $D\left(\dfrac{12}{7};\dfrac{12}{7};0\right)$.\\
		Lại có $\dfrac{ID}{\sin IAD}=\dfrac{AD}{\sin AID}$ và $\dfrac{IO}{\sin IAO}=\dfrac{AO}{\sin AIO}$ và $\widehat{IAD}=\widehat{IAO}$ và $\widehat{AID}+\widehat{AIO}=180^\circ$ nên $\dfrac{ID}{AD}=\dfrac{IO}{AO}$.\\
		\allowdisplaybreaks
		\begin{eqnarray*}
			\vec{AD}=\left(\dfrac{5}{7};-\dfrac{2}{7};2\right)&\Rightarrow& A D=\dfrac{15}{7}.\\
			\vec{ID}=-\dfrac{AD}{AO} \cdot \vec{IO}=-\dfrac{5}{7} \vec{IO} &\Rightarrow& \vec{OI}=\dfrac{7}{12} \vec{OD} \Rightarrow I(1;1;0).
		\end{eqnarray*}
		Do đó $a-b+c=0$.
}\end{ex}
\begin{ex}%[2H2V2-5] 
	Trong không gian với hệ trục tọa độ $Oxyz$, cho tứ diện $ABCD$ có $A(2;-1;1)$, $ B(3;0;-1)$, $C(2;-1;3)$, $D\in O y$ và có thể tích bằng $5$. Tính tổng tung độ của các điểm $D$.
	
	\shortans{-6}
	\loigiai{
		Do $D\in Oy \Rightarrow D(0;y;0)$.\\
		Ta có $\vec{AD}=(-2;y+1;-1)$, $\vec{AB}=(1;1;-2)$, $\vec{AC}=(0;0;2)$.\\
		Khi đó $\left[\vec{AB},\vec{AC}\right]=(2;-2;0)$
		\allowdisplaybreaks
		\begin{eqnarray*}
			V_{ABCD}=\dfrac{1}{6}\left|\left[\vec{AB},\vec{AC}\right]\cdot \vec{AD}\right|=5\Leftrightarrow \left|2(-2)+(-2)(y+1)+0(-1)\right|=30 \Leftrightarrow |-6-2y|=30 \Leftrightarrow \hoac{&y=-18\\&y=12.}
		\end{eqnarray*}
		Vậy $y_1+y_2=12-18=-6$.
}
\end{ex}
\Closesolutionfile{ans}
\indapan{6}{ans/ans-2-C5-OTC-De2-KQ}